--- & Tier 2 --- Network Algebra (title) & 1 & \hyperref[sec:status]{Status of this document} \\
--- & Status of this document & 33 & \hyperref[sec:status]{Status of this document} \\
--- & How to study each node & 88 & \hyperref[sec:howtostudy]{How to study each node} \\
--- & Contents / full section list & 113 & \hyperref[sec:status]{Status of this document}, \S\ref{app:sectionmap} \\
--- & Part 0 --- The astrophysics this tier is for & 193 & \S\ref{part:setting} \\
0.1 & The equilibrium hierarchy & 206 & \S\ref{sec:hierarchy} \\
0.1.1 & Three regimes, stated as claims about fluxes & 208 & \S\ref{sec:regimes} \\
0.1.2 & The historical arc --- and why each step mattered & 287 & \S\ref{sec:history} \\
0.1.3 & What the project bets on this hierarchy & 366 & \S\ref{sec:bet} \\
0.2 & Silicon burning, mechanically & 395 & \S\ref{sec:sibmech} \\
0.2.1 & Why ${\sim}3$\,GK is the threshold & 397 & \S\ref{sec:threegk} \\
0.2.2 & Two clusters, one set of bridges & 446 & \S\ref{sec:twoclusters}, Fig.~\ref{fig:clusters} \\
0.2.3 & The Fe-peak endpoint, and why \Ye{} picks it & 502 & \S\ref{sec:fepeak} \\
0.2.4 & Why the flow is not a ladder --- the trap for intuition & 533 & \S\ref{sec:notladder} \\
0.3 & The \Ye{} chain, derived end to end & 554 & \S\ref{sec:yechain} \\
0.3.1 & $\Mch \propto \Ye^2$ --- the derivation & 559 & \S\ref{sec:mch} \\
0.3.2 & Deleptonisation and the homologous core & 593 & \S\ref{sec:delepton}, Fig.~\ref{fig:yechain} \\
0.3.3 & The observable end & 623 & \S\ref{sec:observable} \\
0.3.4 & The error budget this implies & 649 & \S\ref{sec:budget} \\
0.3.5 & The output end --- and the derivative nobody has assembled & 691 & \S\ref{sec:outputend} \\
0.4 & \textbf{Why conservation is structural, not a penalty term} & 743 & \S\ref{sec:whystructural} \\
0.4.1 & Conservation laws are left null vectors & 745 & \S\ref{sec:leftnull} \\
0.4.2 & The soft-penalty failure mode, quantified & 769 & \S\ref{sec:softpenalty} \\
0.4.3 & The reachable manifold, restated & 819 & \S\ref{sec:reachable} \\
0.5 & What physically breaks when a leak happens & 839 & \S\ref{sec:leak} \\
0.5b & Deployment --- what consumes the output, and what speedup is achievable & 874 & \S\ref{sec:deployment} \\
--- & The cost ceiling --- why the fallback rate matters more than the speed & 914 & \S\ref{sec:costceiling} \\
0.6 & The decision this tier serves & 971 & \S\ref{sec:killtest} \\
0.7 & What is \emph{not} in this tier & 1011 & \S\ref{sec:scope} \\
0.8 & Self-check for Part 0 & 1025 & \S\ref{sec:selfcheck-setting} \\
--- & Part I (S7) --- Stoichiometry and the conservation layer & 1048 & \S\ref{part:conservation} \\
I.1 & $\nu$ as a linear map & 1062 & \S\ref{sec:numap} \\
I.2 & Conservation is the left null space --- with the measurement & 1130 & \S\ref{sec:leftnullspace} \\
I.2.1 & The rank measurement --- and why it is the most informative number in S7 & 1134 & \S\ref{sec:rank} \\
I.2.2 & Why the tests can demand exactly \code{0.0} & 1229 & \S\ref{sec:exactzero} \\
I.2b & The same object in another literature --- CRN theory, deficiency, and the flux cone & 1276 & \S\ref{sec:crn} \\
--- & The vocabulary map & 1295 & \S\ref{sec:crnvocab} \\
--- & Deficiency --- measured, and the honest answer & 1310 & \S\ref{sec:deficiency} \\
--- & What CRN theory \emph{does} give: the flux cone & 1360 & \S\ref{sec:fluxcone} \\
I.3 & The weak sector breaks charge conservation --- on purpose & 1402 & \S\ref{sec:weakbreaks} \\
I.4 & Lepton ledgers --- and the tautology trap & 1441 & \S\ref{sec:ledgers} \\
I.5 & $\beta^+$ and the annihilated positron & 1495 & \S\ref{sec:betaplus} \\
I.6 & The constraint matrix $C$ & 1537 & \S\ref{sec:cmatrix} \\
I.7 & Target A's theorem & 1576 & \S\ref{sec:targetatheorem} \\
I.7b & What Target A does \emph{not} guarantee --- positivity & 1615 & \S\ref{sec:positivity} \\
--- & The geometry & 1623 & \S\ref{sec:posgeometry} \\
--- & Why it is not a small problem here & 1639 & \S\ref{sec:posnotsmall} \\
--- & What the literature offers & 1660 & \S\ref{sec:poslit} \\
--- & The minimum that should exist now & 1689 & \S\ref{sec:posminimum} \\
I.8 & The Target B projector, derived twice & 1698 & \S\ref{sec:projector} \\
--- & Derivation 1 --- minimum correction (Lagrange) & 1702 & \S\ref{sec:projlagrange} \\
--- & Derivation 2 --- QR, which is what ships & 1726 & \S\ref{sec:projqr} \\
I.9 & Why the projector acts on the extended vector --- the laundering trap & 1767 & \S\ref{sec:laundering} \\
I.10 & Target A and Target B reach the same set & 1813 & \S\ref{sec:equivalence} \\
I.10b & The effective dimension of the reachable set --- and where \Ye{} hides & 1862 & \S\ref{sec:effdim} \\
--- & The measurement & 1870 & \S\ref{sec:effdim-measurement} \\
--- & The part that changes a design decision & 1913 & \S\ref{sec:effdim-design} \\
--- & The arithmetic that settles it & 1942 & \S\ref{sec:effdim-arithmetic} \\
--- & And it is not an accident --- it is \S I.2.1 & 1968 & \S\ref{sec:effdim-stoich} \\
I.11 & The bipartite graph, the radius, and $K$ & 2011 & \S\ref{sec:graphK} \\
I.12 & Code walkthrough & 2055 & \S\ref{sec:codes7} \\
I.13 & The oracle: what the 30-test gate actually asserts & 2107 & \S\ref{sec:oracles7} \\
I.14 & Self-check for S7 & 2167 & \S\ref{sec:selfcheck-s7} \\
--- & Part II (S8) --- Fluxes, pairs, and cancellation & 2189 & \S\ref{part:fluxes} \\
II.1 & From $\lambda$ to $R$ to \fp/\fm & 2198 & \S\ref{sec:lambdatoR} \\
II.2 & The pair map & 2242 & \S\ref{sec:pairmap} \\
II.3 & The net identity, proved & 2298 & \S\ref{sec:netidentity} \\
II.4 & \kap, three ways & 2338 & \S\ref{sec:kappa} \\
--- & (a) \kap{} as a detailed-balance diagnostic & 2342 & \S\ref{sec:kappa-db} \\
--- & (b) \kap{} as a condition number & 2363 & \S\ref{sec:kappa-cond} \\
--- & (c) \kap{} as a thermodynamic affinity --- the bridge to S9 & 2391 & \S\ref{sec:kappa-affinity} \\
--- & (d) \ldots and therefore \kap{} as entropy production & 2437 & \S\ref{sec:kappa-entropy} \\
II.4b & Entropy production --- a derived active set, a second-law constraint, and a Lyapunov bound & 2473 & \S\ref{sec:entropy} \\
--- & This derives the active set & 2481 & \S\ref{sec:entropyactive} \\
--- & The second-law constraint on a learned flux & 2527 & \S\ref{sec:secondlaw} \\
--- & The Lyapunov corollary --- rollout stability, structurally & 2590 & \S\ref{sec:lyapunov} \\
II.5 & Weak columns: $\kap \equiv 1$, structurally & 2655 & \S\ref{sec:weakkappa} \\
II.6 & The store: why \fm, \netflux, \kap{} are not stored & 2685 & \S\ref{sec:store} \\
II.7 & The $\dd t = 10^{-6}$ label handshake & 2716 & \S\ref{sec:handshake} \\
II.7b & The conditioning variable this suggests --- $\Delta t/\tau$, not $\log\Delta t$ & 2792 & \S\ref{sec:dttau} \\
--- & The argument & 2800 & \S\ref{sec:dttau} \\
--- & The catch, stated properly & 2825 & \S\ref{sec:dttau} \\
II.8 & Cancellation-aware residuals & 2847 & \S\ref{sec:cancresid} \\
II.9 & The measurement that reframed the project: the vacuous pass & 2899 & \S\ref{sec:vacuous} \\
II.10 & Code and oracles & 2967 & \S\ref{sec:codes8} \\
II.11 & Self-check for S8 & 3001 & \S\ref{sec:selfcheck-s8} \\
--- & Part III (S9) --- NSE and QSE & 3019 & \S\ref{part:equilibrium} \\
III.1 & Equilibrium from the ensemble & 3034 & \S\ref{sec:ensemble} \\
--- & The distribution & 3036 & \S\ref{sec:distribution} \\
--- & The equilibrium condition & 3069 & \S\ref{sec:eqcondition} \\
III.2 & Where the binding energy comes from & 3082 & \S\ref{sec:binding} \\
III.3 & The Saha coefficient, term by term & 3138 & \S\ref{sec:saha} \\
III.4 & The two constraints & 3190 & \S\ref{sec:twoconstraints} \\
III.5 & Newton with an analytic Jacobian & 3218 & \S\ref{sec:newton} \\
III.6 & The monotonicity proof --- why the bisection fallback works & 3246 & \S\ref{sec:monotonicity} \\
III.7 & Batching, and why the scalar path survives & 3333 & \S\ref{sec:batching} \\
III.8 & QSE: the cluster and \uG & 3355 & \S\ref{sec:qsecluster} \\
III.9 & QSE as a slow manifold & 3410 & \S\ref{sec:slowmanifold} \\
III.10 & \textbf{The two departure diagnostics --- and they are not the same} & 3456 & \S\ref{sec:departure} \\
III.11 & The group boundary is a config variant, not a fact & 3520 & \S\ref{sec:groupboundary} \\
III.12 & \textbf{What was measured --- three negative results, read carefully} & 3555 & \S\ref{sec:solververify} (1), \S\ref{sec:negresults} (2), (3) \\
III.13 & Code and oracles & 3613 & \S\ref{sec:codes9} \\
III.14 & Self-check for S9 & 3639 & \S\ref{sec:selfcheck-s9}, \S\ref{sec:selfcheck-s8} (Q7, Q8) \\
--- & Part IV (S10) --- Stiff integration and the augmented flux state & 3660 & \S\ref{part:integration} \\
IV.1 & Stiffness, from the Jacobian & 3668 & \S\ref{sec:stiffness} \\
IV.2 & Why BDF, and what L-stability buys & 3703 & \S\ref{sec:bdf} \\
IV.3 & The augmented state $[Y, \Phifl]$ & 3751 & \S\ref{sec:augmented} \\
IV.4 & Why \Phifl{} cannot be recovered from $\Delta Y$ & 3780 & \S\ref{sec:phinotrecoverable} \\
IV.4b & The flux head is unidentifiable under a $\Delta Y$-only loss & 3815 & \S\ref{sec:unidentifiable} \\
IV.5 & $\partial R/\partial Y$ by the product rule & 3871 & \S\ref{sec:productrule} \\
IV.6 & Conservation by construction, inside the solver & 3942 & \S\ref{sec:solverconservation} \\
IV.7 & Error control for \Phifl & 3970 & \S\ref{sec:errorcontrol} \\
IV.8 & The energy identity is a \emph{data} check, not a dynamics check & 3998 & \S\ref{sec:energyidentity} \\
IV.8b & Wegscheider --- the rate set is not globally consistent, measured & 4064 & \S\ref{sec:wegscheider} \\
--- & The condition & 4070 & \S\ref{sec:wegscheider-condition} \\
--- & The measurement & 4094 & \S\ref{sec:wegscheider-measurement} \\
--- & Reading it & 4123 & \S\ref{sec:wegscheider-reading} \\
--- & Why this matters, quantitatively & 4155 & \S\ref{sec:wegscheider-why} \\
IV.9 & The cost wall, diagnosed & 4209 & \S\ref{sec:costwall} \\
IV.10 & Code and oracles & 4255 & \S\ref{sec:codes10} \\
IV.11 & Self-check for S10 & 4289 & \S\ref{sec:selfcheck-s10} \\
--- & Part V --- Executable policy, and the measurement discipline it does not encode & 4310 & \S\ref{part:policy} \\
V.1 & The invariants encoded in code that refuses & 4312 & \S\ref{sec:refuses}, \S\ref{sec:discipline-intro} \\
V.2 & The sampling measure on the regime box & 4369 & \S\ref{sec:samplingmeasure} \\
--- & The measurement & 4382 & \S\ref{sec:samplingmeasure-measurement} \\
--- & Reading it honestly & 4408 & \S\ref{sec:samplingmeasure-reading} \\
V.3 & Cluster-robust inference --- the effective sample size is ${\sim}300$, not 657,000 & 4455 & \S\ref{sec:clusterrobust} \\
--- & The problem & 4457 & \S\ref{sec:clusterrobust-problem} \\
--- & The measurement & 4487 & \S\ref{sec:clusterrobust-measurement} \\
--- & What this does and does not invalidate & 4521 & \S\ref{sec:clusterrobust-scope} \\
--- & The fix, which is cheap & 4549 & \S\ref{sec:clusterrobust-fix} \\
V.4 & Verification practice --- three standard steps that are absent & 4571 & \S\ref{sec:vandv} \\
--- & Part VI --- The open register & 4620 & \S\ref{part:register} \\
VI.1 & The register & 4637 & \S\ref{sec:register} \\
VI.2 & The two items with no home section & 4667 & \S\ref{sec:nohome} \\
--- & R-24 --- Learning-theory prerequisites & 4669 & \S\ref{sec:r24} \\
--- & R-25 --- Shelf life under library and code drift & 4693 & \S\ref{sec:r25} \\
VI.3 & Ranked & 4720 & \S\ref{sec:ranked} \\
--- & Part VII --- How this register was built, and whether it is finished & 4738 & \S\ref{sec:auditmethod} \\
VII.1 & The method --- two passes, and the axes each swept & 4745 & \S\ref{sec:twopasses} \\
VII.2 & Candidate axes that came back covered & 4777 & \S\ref{sec:coveredaxes} \\
VII.3 & Is the audit converging? & 4799 & \S\ref{sec:converging} \\
--- & Where Tier 2 hands off & 4842 & \hyperref[sec:handoff]{Where Tier 2 hands off} \\
--- & What Tier 2 upgraded & 4844 & \hyperref[sec:upgraded]{What Tier 2 upgraded} \\
--- & Threads that carry forward & 4868 & \hyperref[sec:threads]{Threads that carry forward} \\
--- & The three ideas to carry, if you carry nothing else & 4891 & \hyperref[sec:threeideas]{The three ideas to carry} \\
--- & And one habit, from the negative results & 4910 & \hyperref[sec:habit-negative]{And one habit} \\
--- & And one more, from building the register & 4927 & \hyperref[sec:habit-register]{And one more} \\
--- & References & 4948 & \hyperref[sec:references]{References} \\
